\section*{Chapitre \ref{ch:b2:population-genetics} --- Génétique des populations}

\begin{solution}{exo:b2:population-genetics:1}
$p(M) = (1280 + 320)/2000 = 0.8$, $q(N) = 0.2$. Attendus $MM$~:
$0.64\times 1000 = 640$~; $MN$~: $2\times 0.8\times 0.2\times 1000 =
320$~; $NN$~: 40 --- exactement les effectifs observés. La population est
à l'équilibre pour ce locus.
\end{solution}

\begin{solution}{exo:b2:population-genetics:2}
$q = \sqrt{1/20000} = 0.0071$~; porteurs $2pq \approx 0.014$, une
personne sur 70. Les porteurs sont plus nombreux que les personnes
atteintes dans un rapport de $2p/q \approx
280$ à un.
\end{solution}

\begin{solution}{exo:b2:population-genetics:3}
\omterm{def:b2:population-genetics:fitness}{Valeur sélective}~: le nombre attendu de descendants d'un \omterm{def:b2:meiosis-heredity:vocabulary}{génotype}
rapporté à celui du meilleur. \omterm{def:b2:population-genetics:fitness}{Coefficient de sélection} $s$~: le déficit
de \omterm{def:b2:population-genetics:fitness}{valeur sélective}, $w = 1 -
s$. Héritabilité $h^{2}$~: la fraction de la variance phénotypique qui
est génétique additive, ou, de façon équivalente, la pente de la
régression des descendants sur les parents. Forces~: \omterm{def:b2:genome-diversification:mutation}{mutation},
sélection, dérive, migration, accouplements non aléatoires (cette
dernière modifie les fréquences génotypiques, non les \omterm{def:b2:population-genetics:frequencies}{fréquences
alléliques}).
\end{solution}

\begin{solution}{exo:b2:population-genetics:4}
Directionnelle~: le mélanisme de la phalène du bouleau, l'agent étant
les oiseaux qui chassent à vue sur les troncs noircis par la suie.
Stabilisante~: le poids de naissance humain, l'agent étant la mortalité
des bébés très petits et très gros. Balancée~: la drépanocytose, les
agents étant à la fois le paludisme (contre $AA$) et l'anémie (contre
$SS$).
\end{solution}

\begin{solution}{exo:b2:population-genetics:5}
$q_t = q_0/(1 + tq_0)$~: la division par deux demande $1/q_0 = 50$
générations~; atteindre $0.001$ demande $1/0.001 - 1/0.02 = 950$
générations. Presque toutes les copies de l'\omterm{def:b2:meiosis-heredity:vocabulary}{allèle} se trouvent chez des
\omterm{def:b2:meiosis-heredity:vocabulary}{hétérozygotes}, invisibles pour la sélection~; empêcher les malades de se
reproduire (ce que fait déjà la nature, puisqu'ils meurent) ne change
rien de mesurable à l'échelle d'une vie humaine.
\end{solution}

\begin{solution}{exo:b2:population-genetics:6}
$p = 0.3$~: $\bar w = 0.09 + 0.42 + 0.8\times 0.49 = 0.902$~; $p' =
(0.09 + 0.21)/0.902 = 0.333$~; $\Delta p = +0.033$. $p = 0.9$~: $\bar w
= 0.998$, $p' = 0.9018$, $\Delta p = +0.0018$. À $p = 0.9$, seule une
fraction $q^{2} = 0.01$ de la population est $aa$~: l'\omterm{def:b2:meiosis-heredity:vocabulary}{allèle} contre
lequel agit la sélection se cache chez les \omterm{def:b2:meiosis-heredity:vocabulary}{hétérozygotes}, et
$\Delta p = spq^{2}/\bar w$ contient le facteur $q^{2}$.
\end{solution}

\begin{solution}{exo:b2:population-genetics:7}
Avec $s = 0.12$ contre $AA$ et $t = 0.8$ contre $SS$, $\hat q =
s/(s + t) = 0.12/0.92 = 0.13$. Naissances d'enfants malades
$\hat q^{2} = 0.017$, environ un enfant sur 60 --- le prix de la
protection dont bénéficient les \omterm{def:b2:meiosis-heredity:vocabulary}{hétérozygotes}.
\end{solution}

\begin{solution}{exo:b2:population-genetics:8}
Létal \omterm{def:b2:meiosis-heredity:vocabulary}{récessif}~: $\hat q = \sqrt{\mu} = \sqrt{2\times 10^{-5}} =
0.0045$~; naissances atteintes $\hat q^{2} = \mu = 2\times 10^{-5}$, une
sur \num{50000}. Létal \omterm{def:b2:meiosis-heredity:vocabulary}{dominant}~: chaque copie est éliminée dans la
génération où elle apparaît, si bien que les naissances atteintes sont
simplement les nouvelles \omterm{def:b2:genome-diversification:mutation}{mutations}, $2\mu = 4\times 10^{-5}$ (deux
gamètes par naissance), une sur
\num{25000} --- le \omterm{def:b2:meiosis-heredity:vocabulary}{récessif} cache ses \omterm{def:b2:meiosis-heredity:vocabulary}{allèles} dans un rapport de deux
cents à un, le \omterm{def:b2:meiosis-heredity:vocabulary}{dominant} ne cache rien.
\end{solution}

\begin{solution}{exo:b2:population-genetics:9}
$H_t = 0.5\,(1 - 1/100)^{t}$~: après 10 générations, $0.45$~; après 50,
$0.30$~; après 200, $0.067$. Pour $(1 - 1/2N)^{100} = 0.9$~:
$100/(2N) \approx -\ln 0.9 = 0.105$, donc $N \approx 475$, environ 500
reproducteurs.
\end{solution}

\begin{solution}{exo:b2:population-genetics:10}
Vingt copies de gènes~: $P(\text{absent}) = 0.9^{20} = 0.12$. Une copie
unique parmi 20 a la fréquence $0.05$, et un \omterm{def:b2:meiosis-heredity:vocabulary}{allèle} neutre se fixe avec
une probabilité égale à sa fréquence~: $0.05$. Sur un continent de,
disons, \num{100000} oiseaux, une copie unique a une probabilité de
$5\times 10^{-6}$ de se fixer. L'île est un lieu où les \omterm{def:b2:meiosis-heredity:vocabulary}{allèles} rares
sont perdus et où ceux qui survivent peuvent l'emporter~: la dérive à la
fois appauvrit et différencie.
\end{solution}

\begin{solution}{exo:b2:population-genetics:11}
$R = h^{2}S = 0.3\times 1.4\times \qty{1000}{L} = \qty{420}{L}$ par
génération. Avec des taureaux choisis parmi le \qty{1}{\%} meilleur, la
différentielle moyenne vaut $(1.4 + 2.7)/2 = 2.05$ écarts-types et $R =
0.3\times 2050 = \qty{615}{L}$. La réponse ralentit parce que la
sélection épuise la variance additive --- les \omterm{def:b2:meiosis-heredity:vocabulary}{allèles} favorables se
fixent et $h^{2}$ diminue --- et parce que les modifications corrélées
d'autres caractères (fertilité, santé) imposent des coûts qu'une
sélection sur le seul rendement ne voit pas.
\end{solution}

\begin{solution}{exo:b2:population-genetics:12}
La sélection voit des \omterm{def:b2:meiosis-heredity:vocabulary}{phénotypes}. Un \omterm{def:b2:meiosis-heredity:vocabulary}{allèle} \omterm{def:b2:meiosis-heredity:vocabulary}{récessif} chez un
\omterm{def:b2:meiosis-heredity:vocabulary}{hétérozygote} ne produit aucun \omterm{def:b2:meiosis-heredity:vocabulary}{phénotype}, si bien que la sélection ne
peut pas l'éliminer, et sa fréquence diminue comme $1/t$, de plus en
plus lentement. Une \omterm{def:b2:genome-diversification:mutation}{mutation} neutre ne produit aucune différence de
\omterm{def:b2:population-genetics:fitness}{valeur sélective}~; son sort dépend de la seule dérive, et elle se fixe
avec la probabilité $1/2N$. Un caractère d'héritabilité nulle peut être
fortement sélectionné --- les survivants peuvent différer beaucoup de la
population --- et pourtant la génération suivante est inchangée, parce
que les différences n'étaient pas héritées. La sélection ne peut agir que
sur des différences héritables qui modifient la \omterm{def:b2:population-genetics:fitness}{valeur sélective}~; tout
le reste évolue par hasard, ou n'évolue pas.
\end{solution}

\begin{solution}{pb:b2:population-genetics:1}
\textbf{1.} Fraction mélanique $1 - q^{2} = 1 - 0.999^{2} = 0.002$,
une phalène sur 500~; la fraction des \omterm{def:b2:meiosis-heredity:vocabulary}{allèles} $M$ portés par des
\omterm{def:b2:meiosis-heredity:vocabulary}{hétérozygotes} vaut
$2pq/(2pq + 2p^{2}) = q = 0.999$~: pratiquement tous.
\textbf{2.} Avec $\bar w = 1 - sq^{2}$, les copies de l'\omterm{def:b2:meiosis-heredity:vocabulary}{allèle} clair
après sélection sont $pq\cdot 1 + q^{2}(1 - s)$ sur $\bar w$, donc $q'
= (pq + q^{2} - sq^{2})/(1 - sq^{2}) = q(1 - sq)/(1 - sq^{2})$.
\textbf{3.} $q_1 = 0.999\times(1 - 0.2997)/(1 - 0.2994) = 0.99857$~:
une variation de $-0.0014$.
\textbf{4.} $q = \sqrt{0.1} = 0.32$.
\textbf{5.} La récurrence, itérée, demande 28 générations pour $s =
0.3$~; l'approximation donne le même ordre de grandeur (la variation vaut
$-0.0003$ au début, puis s'accélère à mesure que $p$ augmente). Les 50
générations observées correspondent à un $s$ plus faible, environ $0.2$ (la récurrence
donne 45 générations pour $s = 0.2$)~: une sélection de \qtyrange{20}{30}{\%}
par génération, énorme au regard de la plupart des \omterm{def:b2:meiosis-heredity:vocabulary}{allèles}.
\textbf{6.} Avec $w_{MM} = w_{Mm} = 1 - s$, $w_{mm} = 1$~: $\bar w = 1
- s(1 - q^{2})$ et $p' = p(1 - s)/\bar w$, donc $\Delta p = p[(1 - s) -
\bar w]/\bar w = -spq^{2}/\bar w$. La forme claire n'est favorisée que
chez l'\omterm{def:b2:meiosis-heredity:vocabulary}{homozygote} $mm$, une fraction $q^{2} = 0.1$ au départ~: le
retour commence lentement. Quand $q \to 1$, $\Delta p \approx -sp$ et $M$
décline exponentiellement, d'un facteur $1 - s$ à chaque génération~: un
\omterm{def:b2:meiosis-heredity:vocabulary}{allèle} \omterm{def:b2:meiosis-heredity:vocabulary}{dominant} n'a nulle part où se cacher. L'essor était l'image
inverse~:
$\Delta p \approx +sp$ tant que $M$ était rare (début exponentiel), puis
$\Delta q \approx -sq^{2}p \to$ une lente reptation, l'\omterm{def:b2:meiosis-heredity:vocabulary}{allèle} clair se
cachant chez les \omterm{def:b2:meiosis-heredity:vocabulary}{hétérozygotes}. La récurrence donne 27 générations pour
le retour de $p = 0.68$ à $0.001$, dont 16 pour amener l'\omterm{def:b2:meiosis-heredity:vocabulary}{allèle} à \qty{5}{\%}.
\textbf{7.} $24$ copies de gènes~: $0.95^{24} = 0.29$.
\textbf{8.} $H_1 = 0.4\,(1 - 1/24) = 0.383$.
\textbf{9.} $H_{100} = 0.383\,(1 - 1/80)^{100} = 0.383\times 0.284 =
0.109$~: \qty{27}{\%} du $0.4$ du continent.
\textbf{10.} $P(\text{fixation}) = 1/(2N) = 1/80 = 0.0125$~; temps attendu
jusqu'à la fixation $4N = 160$ générations.
\textbf{11.} Sans consanguinité, $q^{2} = 0.0025$~; avec $F = 0.3$,
$q^{2} + Fpq = 0.0025 + 0.3\times 0.95\times 0.05 = 0.017$~: sept
fois plus d'oiseaux atteints.
\textbf{12.} Un migrant par génération ($m = 1/40$) suffit à empêcher
l'île de diverger aux loci neutres et à reconstituer les \omterm{def:b2:meiosis-heredity:vocabulary}{allèles} que la
dérive a perdus~; il importe aussi les \omterm{def:b2:meiosis-heredity:vocabulary}{allèles} délétères du continent à
leur fréquence continentale, qui, sous l'effet de la consanguinité de
l'île, sont exposés chez des \omterm{def:b2:meiosis-heredity:vocabulary}{homozygotes} --- le fardeau est fixé par
l'équilibre entre immigration et exposition.
\textbf{13.} $\hat q = \sqrt{10^{-5}} = 0.0032$~; naissances atteintes $\hat
q^{2} = 10^{-5}$.
\textbf{14.} $\hat q = \sqrt{10^{-5}/0.1} = 0.01$~; naissances atteintes
$10^{-4}$~: une sélection dix fois plus faible rend l'\omterm{def:b2:meiosis-heredity:vocabulary}{allèle} trois fois
plus commun et la maladie dix fois plus.
\textbf{15.} $\hat p = \mu/(hs) = 10^{-5}/0.05 = 2\times 10^{-4}$~;
porteurs $2\hat p = 4\times 10^{-4}$.
\textbf{16.} De $\hat q = 0.0032$ à $0.01$~: l'\omterm{def:b2:meiosis-heredity:vocabulary}{allèle} augmente à un
rythme fixé par la \omterm{def:b2:genome-diversification:mutation}{mutation}, $\mu = 10^{-5}$ par génération, si bien que
l'approche prend de l'ordre de $\Delta q/\mu \approx 700$ générations ---
\num{20000} ans~: toute décision prise aujourd'hui se fera sentir dans
une population qui ne ressemblera en rien à la nôtre.
\textbf{17.} Chaque locus apporte $2\hat q = 0.0063$ copie létale
par personne~; sur \num{20000} loci, environ $130$ --- mais cela suppose
que tout gène peut muter en un \omterm{def:b2:meiosis-heredity:vocabulary}{récessif} létal~; avec l'estimation
habituelle de quelques milliers de gènes essentiels, une personne porte
de plusieurs à quelques dizaines d'équivalents létaux à l'état
\omterm{def:b2:meiosis-heredity:vocabulary}{hétérozygote}.
\textbf{18.} Chacun se trouve à un locus différent et chacun est rare,
si bien que la probabilité que les deux membres d'un couple pris au
hasard portent le même est faible~; des cousins partagent un huitième de
leurs gènes par ascendance, si bien que, pour chacun des létaux d'un
cousin, le risque que l'enfant soit \omterm{def:b2:meiosis-heredity:vocabulary}{homozygote} est de
$1/16$ par létal partagé --- l'excès de mortalité infantile des mariages
entre cousins est mesuré à quelques pour cent.
\textbf{19.} $S = 1.16\times \qty{800}{L} = \qty{930}{L}$~; $R = 0.4\times
930 = \qty{370}{L}$.
\textbf{20.} Cinq générations~: environ \qty{1860}{L}~; moins en
pratique, parce que la variance additive s'épuise, parce que $h^{2}$
a été estimée dans le troupeau d'origine, et parce que l'environnement
(alimentation, logement) qui a fait les meilleures vaches peut ne pas se
transmettre.
\textbf{21.} Différentielle moyenne $(1.16 + 2.4)/2 = 1.78$ écart-type,
$S = \qty{1424}{L}$, $R = 0.4\times 1424 = \qty{570}{L}$.
\textbf{22.} Des \omterm{def:b2:asexual-reproduction:clone}{clones} n'ont aucune variance génétique, donc $h^{2} = 0$ et $R =
0$ quel que soit $S$~: l'héritabilité est une propriété d'une population,
la part de sa variance qui est génétique, et non une propriété du
caractère.
\textbf{23.} $R = 0.7\times 0.6 = \qty{0.42}{mm}$ de plus.
\textbf{24.} La même réponse de $\qty{370}{L}$ exige $S = 370/0.1 =
\qty{3700}{L}$, soit $4.6$ écarts-types --- en ne gardant qu'une vache
sur plus de dix mille~: irréalisable, et c'est pourquoi les caractères
peu héritables s'améliorent lentement ou par d'autres moyens
(testage sur descendance, familles plus nombreuses).
\textbf{25.} Phalène~: $s \approx 0.2$ à $0.3$, 28 générations pour $0.3$
et 45 pour $0.2$, contre 50 observées~; île~: $H_{100} \approx 0.11$,
\qty{27}{\%} de celle du continent~; troupeau~: \qty{370}{L} par
génération avec les vaches seules, \qty{570}{L} avec des taureaux
sélectionnés.
\end{solution}
